\chapter*{Introduction}
\addcontentsline{toc}{chapter}{Introduction}
\markboth{Introduction}{Introduction}

\section*{Motivation}
An exotic option is not priced in a vacuum. A barrier, an Asian or an autocallable must carry a price that is
consistent with the European options the market quotes every day: those options fix the implied-volatility surface,
and a model that does not reproduce it prices the exotic in a market that does not exist.

The simplest model, Black's, uses a single volatility for all strikes and all expiries, so it does not see the smile.
Stochastic-volatility models such as Heston and SABR do see it, but at a cost:

\begin{itemize}
\item the pricing formulas are no longer in closed form, so numerical methods are needed (Fourier, Monte
Carlo, asymptotic expansions);
\item every numerical method has its own ways of going wrong, often silently;
\item the parameters must be estimated from data, and the estimate carries an error that one must be able to measure.
\end{itemize}

A price is worth only as much as the verification behind it. This work starts from a practical question: can one
build a small, readable pricer in which \emph{every} number is reproducible and checked by at least one independent
method? The answer is the \textsc{qpricer} library, and this document is both its study guide and its technical
report.

\section*{What the library does}
\textsc{qpricer} is written in Python and uses only NumPy, SciPy, pandas and matplotlib (plus yfinance to download the
data). It covers four blocks.

\begin{enumerate}
\item \textbf{Market data.} It downloads a snapshot of the option chain on the S\&P\,500 index (SPX), cleans it and,
for each expiry, extracts the forward $F$ and the discount factor $\DF$ from put--call parity, without assuming any
rates or dividends (Chapter~\ref{cap:dati}).
\item \textbf{Models and European options.} Black-76 with Greeks and implied volatility (Chapter~\ref{cap:black});
Heston with its characteristic function, the COS method, the Lewis formula and Monte Carlo with Andersen's QE scheme
(Chapters~\ref{cap:heston} and~\ref{cap:mc}); SABR with $\beta=1$, Hagan's formula and Monte Carlo
(Chapter~\ref{cap:sabr}).
\item \textbf{Exotic options.} Barriers with continuous and discrete monitoring (Chapter~\ref{cap:barriere}),
arithmetic Asians with a geometric control variate (Chapter~\ref{cap:asiatiche}) and single-underlying
autocallables (Chapter~\ref{cap:autocallable}), under Black and under Heston.
\item \textbf{Calibration and benchmark.} Heston calibrated to the whole SPX surface by nonlinear least squares, SABR
expiry by expiry, no-arbitrage checks (Chapter~\ref{cap:calibrazione}); a systematic comparison between Monte Carlo
and closed forms, with convergence rates and variance reduction (Chapter~\ref{cap:benchmark}).
\end{enumerate}

\section*{Architecture}
The code lives in \code{src/qpricer/} and is split into a few modules, each with a single job.
Figure~\ref{fig:architettura} shows how they use one another; Table~\ref{tab:moduli} lists the files.

\begin{figure}[ht]
\centering
\begin{tikzpicture}[
  mod/.style={draw,rounded corners,align=center,text width=33mm,minimum height=13mm,font=\small},
  >=Stealth, every edge/.style={draw,->,thick,black!70}]
\node[mod,fill=bluidea!10]      (market) at (0,0)      {\code{market.py}\\[1pt]$F$, $\DF$ from parity};
\node[mod,fill=bluidea!10]      (data)   at (4.6,0)    {\code{data.py}\\[1pt]SPX chain, OTM surface};
\node[mod,fill=verdeperche!10]  (models) at (0,-2.4)   {\code{models/}\\[1pt]bsm, heston, sabr};
\node[mod,fill=verdeperche!10]  (iv)     at (4.6,-2.4) {\code{iv.py}\\[1pt]implied volatility};
\node[mod,fill=arancioprog!10]  (prod)   at (0,-4.8)   {\code{products.py}\\[1pt]payoffs, term sheets};
\node[mod,fill=arancioprog!10]  (eng)    at (4.6,-4.8) {\code{engines/}\\[1pt]analytic, fourier, mc, vr,
  exotic\_mc};
\node[mod,fill=violadom!10]     (cal)    at (9.4,-2.4) {\code{calibration.py}\\[1pt]objectives, fit, diagnostics};
\path (market) edge (data)
      (data)   edge (cal)
      (models) edge (iv)
      (iv)     edge (cal)
      (models) edge (eng)
      (prod)   edge (eng)
      (eng)    edge (cal);
\end{tikzpicture}
\caption{The modules of \code{src/qpricer/}. An arrow $A\to B$ means that $B$ uses $A$. The scripts in
\code{scripts/} and the tests in \code{tests/} sit on top of everything and produce every number in this document.}
\label{fig:architettura}
\end{figure}

\begin{table}[ht]
\centering
\small
\caption{The files of \code{src/qpricer/} (2644 lines in total, counted with \code{wc -l}).}
\label{tab:moduli}
\begin{tabularx}{\textwidth}{l r X l}
\toprule
File & Lines & Content & Ch.\\
\midrule
\code{market.py} & 233 & put--call parity regression, $F$ and $\DF$ per expiry, forward curve & \ref{cap:dati}\\
\code{data.py} & 270 & chain loading and cleaning, ACT/365F times, OTM surface, arbitrage filter
& \ref{cap:dati}\\
\code{iv.py} & 160 & Black inversion: Newton with Brent fallback & \ref{cap:black}\\
\code{models/bsm.py} & 127 & Black-76 price and Greeks & \ref{cap:black}\\
\code{models/heston.py} & 113 & parameters, characteristic function, cumulants & \ref{cap:heston}\\
\code{models/sabr.py} & 142 & Hagan's formula, ATM inversion, Monte Carlo & \ref{cap:sabr}\\
\code{products.py} & 80 & vanilla payoffs, autocallable term sheets and cash flows & \ref{cap:autocallable}\\
\code{engines/analytic.py} & 208 & barriers (Reiner--Rubinstein, BGK), geometric Asian, autocallable oracles
& \ref{cap:barriere}--\ref{cap:autocallable}\\
\code{engines/fourier.py} & 149 & self-verifying COS and the Lewis formula & \ref{cap:heston}\\
\code{engines/mc.py} & 187 & Heston simulation (Euler, QE, QE-M), estimator and SE & \ref{cap:mc}\\
\code{engines/vr.py} & 162 & antithetic variates, control variate, Sobol' & \ref{cap:benchmark}\\
\code{engines/exotic\_mc.py} & 230 & Monte Carlo for barriers, Asians and autocallables
& \ref{cap:barriere}--\ref{cap:autocallable}\\
\code{calibration.py} & 582 & objectives, residuals, Heston and SABR fits, error bars, arbitrage checks
& \ref{cap:calibrazione}\\
\bottomrule
\end{tabularx}
\end{table}

The split has a precise purpose. \code{products.py} describes \emph{what} a contract pays and knows nothing about
models; \code{models/} describes the dynamics and knows nothing about contracts; \code{engines/} combines the two.
This way the same payoff can be priced with several engines, and it is precisely the comparison between different
engines that makes verification possible.

\section*{Conventions}
\paragraph{Forward-based kernels.} Every pricing function takes $F$, $\DF$ and $T$, never the rate $r$ or the
dividend yield $q$. The reason is that put--call parity gives $F$ and $\DF$ directly for each expiry, and for SPX the
rates, dividends and stock-borrow cost cannot be observed separately. Everything that matters for a European option
is contained in the forward (Chapter~\ref{cap:opzioni}).

\paragraph{Time, rates, volatilities.} Time is measured in years with the ACT/365F convention, computed in seconds up
to the settlement instant. Rates are continuously compounded. Volatilities are annualized and in decimals: $0.20$
means 20\%, and ``one vol point'' is $0.01$.

\paragraph{Arithmetic.} All computations are in double precision (\code{float64}) and vectorized with NumPy, with no
Python loops over paths or strikes.

\paragraph{Reproducible randomness.} Random numbers come only from an explicit generator, created with
\code{np.random.default\_rng(seed)}. Every Monte Carlo estimate returns the quadruple (price, standard error, number
of paths, seed), so anyone can redo the computation and obtain the same digits.

\paragraph{Tests.} The \code{pytest} suite in \code{tests/} contains 599 tests. The frozen reference values stored in
\code{tests/golden/} come from methods independent of the implementation, and the source is stated in the file
itself.

\section*{Contributions and main results}
This work does not propose new models: Black, Heston and SABR are classics. The contribution is a compact and
\emph{verified} implementation of the whole chain, from raw data to the price of an exotic, with the errors of every
step measured and stated. In particular:

\begin{itemize}
\item a data pipeline that extracts $F$ and $\DF$ from parity with a robust regression, and corrects the discount
factors of the short expiries, which the data determine poorly (Chapter~\ref{cap:dati});
\item a \emph{self-verifying} COS method, which chooses its own numerical parameters and checks its own truncation
error, compared against the Lewis formula (Chapter~\ref{cap:heston});
\item at least one independent oracle for every exotic: closed forms, model-free identities, limiting cases in
which Heston reduces to Black (Chapters~\ref{cap:barriere}--\ref{cap:autocallable});
\item a Heston calibration to SPX validated first on synthetic data, with error bars and a sensitivity analysis that
separates estimation error from model error (Chapter~\ref{cap:calibrazione});
\item a Monte Carlo benchmark that measures convergence rates, variance reduction and work-normalized efficiency
(Chapter~\ref{cap:benchmark});
\item an independent validation of every phase, with its own reference computations rather than the
implementation's tests (Chapter~\ref{cap:validazione}).
\end{itemize}

Table~\ref{tab:risultati-intro} collects a few measured numbers; the full picture is in
Chapter~\ref{cap:validazione}.

\begin{table}[ht]
\centering
\small
\caption{Some measured results. ``$|z|$'' is the distance between the Monte Carlo estimate and the reference in
standard errors; the test passes if $|z|\le3$.}
\label{tab:risultati-intro}
\begin{tabularx}{\textwidth}{X l l}
\toprule
Check & Measure & Ch.\\
\midrule
put--call parity in Black-76, $K/F\in[0.25,4]$ & $1.1\times10^{-13}$ & \ref{cap:black}\\
implied-volatility round trip & $\le8.2\times10^{-13}$ & \ref{cap:black}\\
COS vs Lewis on a grid of strikes and expiries & $\le1.8\times10^{-10}$ & \ref{cap:heston}\\
QE and QE-M vs COS (4 parameter sets, 3 strikes) & 24 of 24 within 3 SE & \ref{cap:mc}\\
Hagan vs Monte Carlo for SABR, $\nu^2T\le0.16$ (21 comparisons) & $\max|z|=1.68$ & \ref{cap:sabr}\\
barriers: knock-in $+$ knock-out $=$ vanilla on the same paths & $3.6\times10^{-15}$ & \ref{cap:barriere}\\
barriers: Monte Carlo with Brownian bridge vs Reiner--Rubinstein & $\max|z|=1.99$ & \ref{cap:barriere}\\
geometric Asian: Monte Carlo vs Kemna--Vorst (39 cases) & $\max|z|=1.90$ & \ref{cap:asiatiche}\\
arithmetic Asian: variance reduction with the geometric control & 612--2202 & \ref{cap:asiatiche}\\
quarterly autocallable: calibrated Heston minus Black (Heston ATM vol) & $-0.0187$ (SE 0.0002)
& \ref{cap:autocallable}\\
Heston calibration: starts reaching the same minimum & 16 of 16 & \ref{cap:calibrazione}\\
Heston RMSE in sample / out of sample (vol points) & 0.328 / 1.379 & \ref{cap:calibrazione}\\
log-log slopes of the Monte Carlo error (12 estimates) & from $-0.554$ to $-0.487$ & \ref{cap:benchmark}\\
variance reduction at $2^{20}$ paths, ATM Black: antithetic / control / Sobol'
& 1.6 / 4.7 / $1.4\times10^5$ & \ref{cap:benchmark}\\
\bottomrule
\end{tabularx}
\end{table}

Equally important are the limitations found, which the document reports without softening them. Heston, with a
single time scale, cannot reproduce the skew of the weekly expiries and that of the long ones at the same time: out
of sample the error is about four times the in-sample error. Hagan's formula for SABR is accurate only for small
$\nu^2T$, while the 39 calibrated SPX expiries all lie beyond the verified threshold $\nu^2T\le0.16$. On real data
the uncertainty of the Heston parameters is about ten times what the curvature standard errors suggest.

\section*{How to read this document}
\paragraph{Who it is for.} Readers who know calculus (derivatives, integrals, Taylor series), some probability
(random variable, expected value, variance, normal distribution) and a little linear algebra. Everything else
(Brownian motion, Itô calculus, changes of measure, Fourier transforms, Monte Carlo, nonlinear least squares) is
introduced when needed, intuitively but correctly. Readers who already know the subject can read the Abstract,
Table~\ref{tab:risultati-intro} and Chapter~\ref{cap:validazione}, and dip into the chapters only where needed.

\paragraph{The structure.} The document has four parts.

\begin{itemize}
\item \textbf{Part I, Foundations} (Chapters~\ref{cap:opzioni}--\ref{cap:black}): options, forwards and absence of
arbitrage; the mathematical tools; the SPX data; the Black model.
\item \textbf{Part II, Stochastic volatility and Monte Carlo} (Chapters~\ref{cap:heston}--\ref{cap:sabr}): the
smile and Heston with Fourier methods; Monte Carlo and Heston simulation; SABR.
\item \textbf{Part III, Exotic options} (Chapters~\ref{cap:barriere}--\ref{cap:autocallable}): barriers, Asians,
autocallables.
\item \textbf{Part IV, Calibration, benchmark and validation}
(Chapters~\ref{cap:calibrazione}--\ref{cap:validazione}): calibration to SPX, the Monte Carlo benchmark, a summary of
the validation.
\end{itemize}

The appendices contain a glossary and the bibliography.

\paragraph{A path for readers starting from scratch.} Figure~\ref{fig:percorso} shows the dependencies between the
chapters.

\begin{enumerate}
\item Read Chapters~\ref{cap:opzioni}--\ref{cap:black} in full: they are the basis for everything else. On a first
pass, the first sections and the summary of Chapter~\ref{cap:dati} are enough.
\item Move on to Chapter~\ref{cap:heston} (Heston) and Chapter~\ref{cap:mc} (Monte Carlo). The latter is needed for
all the exotics.
\item Read the exotics in the order \ref{cap:barriere}, \ref{cap:asiatiche}, \ref{cap:autocallable}: each one uses
something from the previous one.
\item Chapter~\ref{cap:sabr} (SABR) can be postponed: it is needed mainly for the calibration.
\item Finish with calibration, benchmark and validation (Chapters~\ref{cap:calibrazione}--\ref{cap:validazione}),
which use almost everything.
\end{enumerate}

At the end of each chapter, the summary and the self-check questions tell you whether you are ready to move on.

\begin{figure}[ht]
\centering
\begin{tikzpicture}[node distance=6mm and 7mm,
  box/.style={draw,rounded corners,align=center,minimum width=27mm,minimum height=9mm,font=\small},
  >=Stealth]
\node[box,fill=bluidea!10] (c1) {Ch.\,\ref{cap:opzioni}--\ref{cap:strumenti}\\options, tools};
\node[box,fill=bluidea!10,right=of c1] (c3) {Ch.\,\ref{cap:dati}\\SPX data};
\node[box,fill=bluidea!10,right=of c3] (c4) {Ch.\,\ref{cap:black}\\Black-76};
\node[box,fill=verdeperche!10,below=of c4] (c5) {Ch.\,\ref{cap:heston}\\Heston (Fourier)};
\node[box,fill=verdeperche!10,left=of c5] (c6) {Ch.\,\ref{cap:mc}\\Monte Carlo};
\node[box,fill=verdeperche!10,left=of c6] (c7) {Ch.\,\ref{cap:sabr}\\SABR};
\node[box,fill=arancioprog!10,below=of c7] (c8) {Ch.\,\ref{cap:barriere}\\barriers};
\node[box,fill=arancioprog!10,right=of c8] (c9) {Ch.\,\ref{cap:asiatiche}\\Asians};
\node[box,fill=arancioprog!10,right=of c9] (c10) {Ch.\,\ref{cap:autocallable}\\autocallables};
\node[box,fill=violadom!10,below=of c8] (c11) {Ch.\,\ref{cap:calibrazione}\\calibration};
\node[box,fill=violadom!10,right=of c11] (c12) {Ch.\,\ref{cap:benchmark}\\MC benchmark};
\node[box,fill=violadom!10,right=of c12] (c13) {Ch.\,\ref{cap:validazione}\\validation};
\draw[->] (c1)--(c3); \draw[->] (c3)--(c4); \draw[->] (c4)--(c5);
\draw[->] (c5)--(c6); \draw[->] (c6)--(c7);
\draw[->] (c6)--(c8); \draw[->] (c6)--(c9); \draw[->] (c6)--(c10);
\draw[->] (c8)--(c9); \draw[->] (c9)--(c10);
\draw[->] (c7.west) -- ++(-5mm,0) |- (c11.west); \draw[->] (c11)--(c12); \draw[->] (c12)--(c13);
\draw[->] (c10)--(c12);
\end{tikzpicture}
\caption{Dependencies between the chapters. Blue: foundations; green: stochastic volatility and Monte Carlo; orange:
exotics; purple: calibration, benchmark and validation.}
\label{fig:percorso}
\end{figure}

\section*{The boxes}
The main text contains the definitions and the derivations. Seven kinds of box accompany it, each with a fixed role,
so that you can tell at a glance what you are reading.

\begin{idea}
The concept to remember even after the formulas have been forgotten.
\end{idea}
\begin{perche}
The reasoning that justifies a choice or a step: why it is done this way and not another.
\end{perche}
\begin{attenzione}
A common mistake, conceptual or numerical, and how to avoid it.
\end{attenzione}
\begin{esempio}
A concrete computation with numbers, often taken from the project.
\end{esempio}
\begin{progetto}
How the concept is implemented in the code (functions and files) and which results were measured. All the numbers
in these boxes come from scripts in the repository and can be regenerated.
\end{progetto}
\begin{riassunto}
The essential points of the chapter, at the end of each one.
\end{riassunto}
\begin{domande}
Questions to check your understanding, with a hint for the answer.
\end{domande}

\section*{Main notation}
\begin{center}
\small
\begin{tabularx}{\textwidth}{>{$}l<{$}X}
\toprule
\text{Symbol} & Meaning \\
\midrule
S_t & \emph{spot} price of the underlying at time $t$ (here the SPX index) \\
F,\ F_t=F(t,T) & \emph{forward} price for delivery at $T$, observed today ($F$) or at time $t$ ($F_t$) \\
\DF=P(0,T) & discount factor: today's value of 1 dollar paid at $T$ \\
K & strike (exercise price) \\
T & expiry in years (ACT/365F: calendar days divided by 365) \\
\sigma & annualized volatility, in decimals ($0.20 = 20\%$);\newline one vol point is $0.01$ \\
x=\ln(F/K),\ k=\ln(K/F) & log-moneyness (two opposite signs; both are used) \\
s=\sigma\sqrt T & total volatility (standard deviation of $\ln F_T$) \\
\theta=\pm1 & $+1$ call, $-1$ put (in the Heston chapter $\theta$ is also a parameter: the context makes it clear) \\
\Phi,\ \varphi & cumulative distribution function and density of the standard normal \\
\Q^T & $T$-forward measure (Chapter~\ref{cap:strumenti}) \\
\E^T[\cdot] & expected value under $\Q^T$ \\
\SE & standard error of a Monte Carlo estimate \\
\varepsilon=2^{-52} & machine epsilon in double precision ($\approx2.2\times10^{-16}$) \\
\bottomrule
\end{tabularx}
\end{center}
